% ---- pg-541 （书页 529） ----

其中 $0<\tilde{s}_{ik}\le 1$ 为第 $i$ 个项目与第 $k$ 个项目之间的相似性, 而 $d_{ik}$ 则为相应的距离.

然而距离必须满足\hyperref[eq:1-25]{式(1-25)}, 所以不是总能由相似性来构造的. 正如高尔(Gower)所证明的那样(见文献[11]和[12]), 仅当相似性矩阵为非负定时才能做到这点. 在非负定条件下以及当最大相似性尺度已标准化(即 $\tilde{s}_{ii}=1$)时, 可以定义
\phantomsection\label{eq:12-9}
\begin{equation}
d_{ik}=\sqrt{2(1-\tilde{s}_{ik})}
\tag{12-9}
\end{equation}
且 $d_{ik}$ 满足距离的性质.

\mvsub{变量对的相似性及关联性量度}

迄今我们一直在对项目讨论相似性量度. 但在某些应用中, 必须进行分组的是变量而不是项目. 变量之间的相似性量度常取样本相关系数的形式. 不仅如此, 在某些聚类应用中, 负相关系数要用它们的绝对值来代替.

当变量为二值变量时, 又可将数据表示为列联表的形式. 不过这一次描述类别的是变量, 而不是项目. 对每一对变量, 有 $n$ 个项目被交叉分类成一个表. 对于通常的 0—1 代码, 此表具有如下形式:

\phantomsection\label{eq:12-10}
\begin{equation}
\begin{array}{ll|cc|c}
\hline
 & & \multicolumn{2}{c|}{变量 $k$} & \\
\cline{3-4}
 & & 1 & 0 & 合\quad 计\\
\hline
变量 $i$ & 1 & $a$ & $b$ & $a+b$\\
 & 0 & $c$ & $d$ & $c+d$\\
\hline
合\quad 计 & & $a+c$ & $b+d$ & $n=a+b+c+d$\\
\hline
\end{array}
\tag{12-10}
\end{equation}

例如, 在 $n$ 个项目中, 变量 $i=1$ 和变量 $k=0$ 的项目有 $b$ 个.

将通常的积矩相关系数公式用于列联表(12-10)中的二值变量, 就得到(见练习 12.3)
\phantomsection\label{eq:12-11}
\begin{equation}
r=\frac{ad-bc}{[(a+b)(c+d)(a+c)(b+d)]^{1/2}}
\tag{12-11}
\end{equation}
这个数可作为两个变量之间相似性的量度.

\hyperref[eq:12-11]{式(12-11)}中的相关系数与检验两个类别变量独立性的 $\chi^2$ 统计量有关($r=\chi^2/n$). 对给定的 $n$, 大的相似性(或相关性)与独立性不足相一致.

给定(12-10)中的表, 便可得出与\hyperref[tab:12.1]{表 12.1} 中所列内容完全类似的关联性(或相似性)量度. 唯一要做的改变是用 $n$(项目数)来代替 $p$(变量数).

\mvsub{相似性问题的小结}

为了对本节进行总结, 我们注意到有很多种方法可用来测量各对对象之间的相似性. 看来, 大多数实际工作者用距离[见\hyperref[eq:12-1]{式(12-1)}到\hyperref[eq:12-5]{式(12-5)}]或\hyperref[tab:12.2]{表 12.2} 中的相似性系数来对项目进行聚类, 而用相关系数对变量进行聚类. 然而, 有时输入聚类算法程序的也许仅仅是简单的频数.

\noindent\textbf{例 12.2}\quad (测量 11 种语言的相似性)

词的含义随着历史进程而变化. 然而数字 $1,2,3,\cdots$ 的意义却代表一个显而易见的例外. 这样, 语言的第一次比较也许可以用数词作比较对象. \hyperref[tab:12.2]{表 12.2} 给出英语、波兰语、匈牙利语及其他八种现代欧洲语言的前 10 个数字. (只考虑使用罗马字母的语言, 而且重音符、变音符、分音符等均予省略.) 对表中数词的拼法作粗略考察后便可发现, 前 5 种语言(英语、挪威语、
